%%%%%%%%%%%%%%%%%%%%%%%%%%%%%%%%%%%%%%%%
\chapter{Resumen}
%%%%%%%%%%%%%%%%%%%%%%%%%%%%%%%%%%%%%%%%

En este trabajo se realizó una revisión sistemática de literatura sobre el reconocimiento de imágenes de polinizadores mediante aprendizaje profundo, con énfasis en arquitecturas, métricas de evaluación y factores técnicos asociados al rendimiento. La búsqueda se efectuó en Scopus mediante una estrategia estructurada, siguiendo el enfoque PRISMA para la identificación, selección y depuración de los estudios, junto con criterios de inclusión y exclusión previamente definidos. El corpus final estuvo conformado por 53 estudios publicados desde 2020, centrados en tareas de clasificación, detección y segmentación. Los resultados evidencian un predominio de la detección y una presencia destacada de la familia YOLO, junto con arquitecturas como ResNet, MobileNet, EfficientNet, Faster R-CNN y Mask R-CNN. Las métricas más reportadas fueron exactitud, precisión y sensibilidad, mientras que en detección destacaron mAP, IoU y FPS. Además, se identificó que el rendimiento depende de la arquitectura y de estrategias como aumento de datos, aprendizaje por transferencia y métodos de optimización. La evidencia revisada muestra avances relevantes en el reconocimiento automatizado de polinizadores e indica que la comparación por tipo de tarea, modelo y métrica permite ordenar una literatura metodológicamente heterogénea.


\paragraph{Palabras clave:} Aprendizaje profundo, Métricas de desempeño, Polinizadores, Reconocimiento de imágenes, Visión por computador.

%%%%%%%%%%%%%%%%%%%%%%%%%%%%%%%%%%%%%%%%
\chapter{Abstract}
%%%%%%%%%%%%%%%%%%%%%%%%%%%%%%%%%%%%%%%%

This study conducted a systematic literature review on pollinator image recognition using deep learning, with emphasis on architectures, evaluation metrics, and technical factors associated with model performance. The search was carried out in Scopus through a structured search strategy, following the PRISMA approach for the identification, screening, and selection of studies, together with previously defined inclusion and exclusion criteria. The final corpus consisted of 53 studies published since 2020, focused on classification, detection, and segmentation tasks. The results show a predominance of detection tasks and a notable presence of the YOLO family, along with architectures such as ResNet, MobileNet, EfficientNet, Faster R-CNN, and Mask R-CNN. The most frequently reported metrics were accuracy, precision, and recall, while mAP, IoU, and FPS were prominent in detection tasks. In addition, model performance was found to depend on the architecture and on strategies such as data augmentation, transfer learning, and optimization methods. The reviewed evidence shows relevant advances in the automated recognition of pollinators and indicates that comparison by task type, model, and metric makes it possible to organize a methodologically heterogeneous body of literature.


\paragraph{Keywords:} Computer vision, Deep learning, Image recognition, Performance metrics, Pollinators.